\chapter{La volatilité comme quantité négociée~: premier contact}\label{ch:m1:volatility-first-contact}

Le call à trois mois sur 100 est offert au cours vendeur de 4.06. Personne sur un desk d'options ne le
dit ainsi. On dit que le call vaut «~vingt vol~»~; que le put 90 vaut «~vingt-quatre~»~;
que le «~skew~» de l'action est de «~six points~»~; que «~le forward est à 100.20, donc le
marché a un dividende de 80 cents dedans~». Les prix en dollars changent à chaque fois que
l'action bouge d'un tick~; ces nombres-là ne changent pas, et ce sont donc ce que le desk cote,
compare, retient et négocie. Ce chapitre apprend ce langage. Il a besoin
d'une formule venue de l'extérieur (le prix de Black--Scholes, dérivé dans le
livre~5 et utilisé ici comme une boîte noire qui transforme une volatilité en prix)
et d'une identité qui ne suppose aucun modèle.

\section{Parité call-put, avec dividendes et emprunt}

\begin{proposition}[Parité call-put]\label{prop:m1:volatility-first-contact:parity}
Pour des options européennes de même strike $K$ et même échéance $T$ sur le même
sous-jacent,
\[
C - P \;=\; \mathrm{e}^{-rT}\,(F - K),
\]
où $F$ est le prix forward du sous-jacent pour $T$~: la \emph{parité
call-put}\index{parité call-put}. Pour une action, $F\,\mathrm{e}^{-rT} = S - D$,
avec $D$ la valeur actuelle des dividendes versés avant $T$, à laquelle s'ajoutent
pour le vendeur à découvert, les \omterm{def:m1:financing:seclending}{frais d'emprunt}.
\end{proposition}
\begin{proof}
Être long d'un call et short d'un put de strike $K$, c'est s'engager à acheter à $K$ en $T$~: un
contrat forward de strike $K$, valant $\mathrm{e}^{-rT}(F-K)$ aujourd'hui.
\end{proof}

La parité tient quelle que soit la volatilité, la distribution ou le modèle. Elle est
imposée par deux opérations qu'un \omterm{def:m1:what-a-trading-firm-does:market-maker}{teneur de marché} peut exécuter mécaniquement~: la
\emph{conversion} (acheter l'action, acheter le put, vendre le call) et le
\emph{reversal} (l'opération inverse, qui suppose un emprunt de titres). Lue à rebours,
elle devient un instrument de mesure~:

\begin{method}[Lire le forward sur la chaîne]\label{met:m1:volatility-first-contact:forward}
\begin{enumerate}
\item Pour chaque strike, calculer $F_K = K + \mathrm{e}^{rT}(C_K - P_K)$ à partir
des prix milieu. Les valeurs doivent concorder à la largeur des cotations près~; faire la
moyenne de celles qui sont au voisinage de la monnaie.
\item Le carry implicite est $D^{\text{imp}} = S - \mathrm{e}^{-rT}F$~: dividendes attendus
plus les \omterm{def:m1:financing:seclending}{frais d'emprunt} que paie l'arbitragiste marginal.
\item Si $D^{\text{imp}}$ dépasse tout dividende plausible, l'action est difficile
à emprunter (voir au \cref{ch:m1:stock-loan-and-short-selling}) et les options signalent
le niveau des frais.
\end{enumerate}
\end{method}

Pour les options américaines, la parité devient une inégalité, puisque les deux côtés peuvent être
exercés par anticipation~; au voisinage de la monnaie et en l'absence de dividendes, la relation européenne
en est une bonne approximation, et elle échoue précisément dans les cas du
\cref{ch:m1:option-contracts-and-exchanges} où l'exercice anticipé est
rationnel.

\section{Volatilité implicite}

\begin{definition}[Volatilité implicite]\label{def:m1:volatility-first-contact:iv}
La \emph{volatilité implicite}\index{volatilité implicite} d'une option est la
valeur du paramètre de volatilité qui, introduit dans la formule de Black--Scholes
avec le strike, l'échéance, le forward et le taux d'intérêt de l'option, redonne
le prix de marché de celle-ci. C'est un changement d'unité, des dollars vers l'écart-type annualisé
des rendements, et il n'affirme rien sur la vérité du modèle.
\end{definition}

\begin{definition}[À la monnaie]\label{def:m1:volatility-first-contact:atm}
Une option est \emph{à la monnaie}\index{a la monnaie@à la monnaie} lorsque son strike est égal
au forward (ou, de manière plus lâche, au spot). Son prix est, à bonne approximation,
proportionnel à la volatilité~:
$C_{\text{ATM}} \approx 0.4\,F\,\mathrm{e}^{-rT}\sigma\sqrt{T}$.
\end{definition}

Le prix d'une option croît avec la volatilité, depuis sa \omterm{def:m1:option-contracts-and-exchanges:series}{valeur intrinsèque}
actualisée en $\sigma = 0$ jusqu'au forward actualisé (call) ou au strike (put) lorsque
$\sigma \to \infty$. Entre ces bornes, il existe exactement une \omterm{def:m1:volatility-first-contact:iv}{volatilité implicite},
que la dichotomie permet de trouver~; en dehors d'elles, il n'en existe aucune, et un
solveur doit le signaler~: un prix inférieur à la \omterm{def:m1:option-contracts-and-exchanges:series}{valeur intrinsèque} traduit un cours périmé
ou un mauvais forward, et non une volatilité négative.

\begin{omfigure}
\begin{tikzpicture}
\begin{axis}[width=10.5cm,height=4.8cm,
  xlabel={volatilité (\% par an)},ylabel={prix de l'option},
  tick label style={font=\small},label style={font=\small},
  xmin=0,xmax=60,ymin=0,ymax=13,grid=major,grid style={gray!25},
  legend columns=2,legend style={font=\footnotesize,at={(0.5,-0.32)},anchor=north,draw=none,fill=none,/tikz/every even column/.append style={column sep=8pt}},
  table/col sep=comma]
\addplot[omDef,very thick] table[x=vol_pct,y=call_100]{figdata/markets-1/25-volatility-first-contact/price_vs_vol.csv};
\addplot[omThm,very thick,dashed] table[x=vol_pct,y=put_90]{figdata/markets-1/25-volatility-first-contact/price_vs_vol.csv};
\draw[thin] (axis cs:0,4.06) -- (axis cs:20.08,4.06) -- (axis cs:20.08,0);
\draw[thin] (axis cs:0,1.12) -- (axis cs:23.82,1.12) -- (axis cs:23.82,0);
\legend{call 100,put 90}
\end{axis}
\end{tikzpicture}

\omcaption{Prix en fonction de la volatilité pour deux options à trois mois sur une action valant
100 avec un forward de 100.20. L'inversion de la courbe au prix de marché donne
la \omterm{def:m1:volatility-first-contact:iv}{volatilité implicite}~: 4.06 correspond à 20.1 pour le call, 1.12 à 23.8 pour le put.
Au voisinage de la monnaie, la courbe est une droite~; loin de celle-ci, elle est plate aux faibles
volatilités, et c'est là que les solveurs perdent en précision. Données~: la
construction du chapitre.}\label{fig:m1:volatility-first-contact:inversion}
\end{omfigure}

\begin{definition}[Volatilité réalisée et variance]\label{def:m1:volatility-first-contact:realised}
La \emph{volatilité réalisée}\index{volatilité réalisée} d'un actif sur une
période est l'écart-type annualisé de ses rendements sur cette
période, calculé par convention à partir des rendements logarithmiques quotidiens, sans soustraire
la moyenne. La \emph{variance}\index{variance} en est le carré. La \omterm{def:m1:volatility-first-contact:iv}{volatilité implicite} est
un prix fixé aujourd'hui~; la volatilité réalisée est une statistique connue après coup~; la
différence entre les deux est ce que gagne ou perd une position d'option couverte.
\end{definition}

\section{Skew et structure par terme}

\begin{definition}[Skew de volatilité et structure par terme]\label{def:m1:volatility-first-contact:skew}
Le \emph{skew de volatilité}\index{skew de volatilité} (ou smile) est la
dépendance de la \omterm{def:m1:volatility-first-contact:iv}{volatilité implicite} au prix d'exercice, pour une échéance donnée~; la \emph{structure
par terme de la volatilité}\index{structure par terme de la volatilité} est sa
dépendance à l'échéance pour une moneyness donnée~; ensemble, ils forment la
surface de volatilité.
\end{definition}

Si le modèle qui sous-tend le changement d'unités était vrai, chaque prix d'exercice aurait
la même \omterm{def:m1:volatility-first-contact:iv}{volatilité implicite}. Pour les indices actions depuis 1987, la \omterm{def:m1:volatility-first-contact:iv}{volatilité implicite}
des bas prix d'exercice dépasse nettement celle des hauts~: les krachs
sont plus fréquents et plus redoutés que ne le permet le modèle, et les investisseurs paient
pour se protéger. Le skew se cote comme un écart en points de volatilité
entre deux prix d'exercice de référence ou deux deltas, et se négocie comme tel
(\cref{fig:m1:volatility-first-contact:smile}).

\begin{omfigure}
\begin{tikzpicture}
\begin{axis}[width=10.5cm,height=4.8cm,
  xlabel={prix d'exercice},ylabel={volatilité implicite (\%)},
  tick label style={font=\small},label style={font=\small},
  xmin=78,xmax=122,ymin=14,ymax=33,grid=major,grid style={gray!25},
  nodes near coords,nodes near coords style={font=\scriptsize,above=2pt,/pgf/number format/.cd,fixed,fixed zerofill,precision=1},
  point meta=explicit,table/col sep=comma]
\addplot[omDef,very thick,mark=*] table[x=strike,y=implied_vol_pct,meta=implied_vol_pct]{figdata/markets-1/25-volatility-first-contact/chain.csv};
\end{axis}
\end{tikzpicture}

\omcaption{Le smile reconstitué à partir des cours affichés à l'écran de la
chaîne synthétique du chapitre, option hors de la monnaie à chaque prix d'exercice, avec le forward
lu à partir de la \omterm{prop:m1:volatility-first-contact:parity}{parité call-put}. Le skew 90--110 vaut 6.2 points de volatilité. Données~: le
tutoriel.}\label{fig:m1:volatility-first-contact:smile}
\end{omfigure}

\section{L'indice de volatilité}

\begin{definition}[Indice de volatilité]\label{def:m1:volatility-first-contact:vix}
Un \emph{indice de volatilité}\index{indice de volatilité} est un nombre publié qui
synthétise la \omterm{def:m1:volatility-first-contact:iv}{volatilité implicite} des options sur un indice pour un horizon fixe,
calculé à partir des prix des options au moyen d'une formule n'utilisant aucun modèle de valorisation~:
\[
\sigma^2 \;=\; \frac{2}{T}\sum_i \frac{\Delta K_i}{K_i^2}\,\mathrm{e}^{RT}\,Q(K_i)
\;-\; \frac{1}{T}\Bigl(\frac{F}{K_0}-1\Bigr)^2 ,
\]
avec $Q(K_i)$ le prix milieu de l'option hors de la monnaie au prix d'exercice $K_i$,
$K_0$ le premier prix d'exercice inférieur ou égal au forward $F$, et l'indice égal à
$100\,\sigma$.
\end{definition}

\begin{dated}{2026-09}{L'indice et ses contrats à terme}\label{dat:m1:volatility-first-contact:vix}
\emph{L'indice.} L'\omterm{def:m1:volatility-first-contact:vix}{indice de volatilité} le plus connu utilise les options sur le S\&P 500 de
deux échéances encadrant 30 jours, et en interpole les \omterm{def:m1:volatility-first-contact:realised}{variances} à un horizon constant
de 30 jours. Seules entrent les options ayant un cours acheteur non nul, au prix milieu~; en
s'éloignant de la monnaie, le strip s'arrête dès que deux prix d'exercice consécutifs
n'ont plus de cours acheteur. $K_0$ est le prix d'exercice égal ou immédiatement inférieur au
forward, lui-même déterminé à partir du prix d'exercice où les prix du call et du put sont les plus proches. \emph{Les contrats à terme.}
Les contrats à terme sur l'indice ont un multiplicateur de 1\,000~\$ et un pas de cotation de 0.05 point,
soit 50~\$. Ils se règlent le mercredi précédant de trente jours le troisième vendredi du
mois suivant, contre un \omterm{def:m1:basis-roll-and-delivery:soq}{cours spécial d'ouverture} (SOQ) de l'indice, calculé à partir des
prix d'ouverture des options sur le S\&P 500 ce matin-là.
\end{dated}

La somme pondère chaque option par $1/K^2$~: les bas prix d'exercice comptent davantage, donc
l'indice monte avec le skew autant qu'avec le niveau général de la \omterm{def:m1:volatility-first-contact:iv}{volatilité
implicite}. Sur la chaîne synthétique de ce chapitre, la volatilité \omterm{def:m1:volatility-first-contact:atm}{à la monnaie}
vaut 20.1, et le chiffre calculé selon la formule de l'indice 23.0
(\cref{fig:m1:volatility-first-contact:weights}). La formule correspond au prix
du portefeuille d'options qui réplique la \emph{\omterm{def:m1:volatility-first-contact:realised}{variance}} que l'indice
réalisera, résultat démontré dans le livre~5~; ici il importe que l'indice soit le prix
d'un portefeuille négocié, et non une opinion.

\begin{omfigure}
\begin{tikzpicture}
\begin{axis}[width=10cm,height=4.4cm,ybar,bar width=6pt,
  xlabel={prix d'exercice (forward à 100)},ylabel={poids (\%)},
  tick label style={font=\small},label style={font=\small},
  xmin=55,xmax=145,ymin=0,ymax=3.4,ymajorgrids,grid style={gray!25},table/col sep=comma]
\addplot[fill=omDef!60,draw=omDef] table[x=strike,y=weight]{figdata/markets-1/25-volatility-first-contact/weights.csv};
\end{axis}
\end{tikzpicture}

\omcaption{Le poids $\Delta K/K^2$ attribué à un dollar de \omterm{def:m1:option-contracts-and-exchanges:option}{prime de l'option} à
chaque prix d'exercice, normalisé. Un dollar de prime sur le put 60 compte cinq fois
autant qu'un dollar sur le call 140. Données~: la
définition.}\label{fig:m1:volatility-first-contact:weights}
\end{omfigure}

\section{Contrats à terme sur la volatilité}

\begin{definition}[Contrat à terme sur l'indice de volatilité]\label{def:m1:volatility-first-contact:vixfut}
Un \emph{contrat à terme sur le VIX}\index{contrat à terme sur le VIX} est un contrat à terme \omterm{def:m1:basis-roll-and-delivery:delivery}{réglé en espèces} sur la valeur
de l'\omterm{def:m1:volatility-first-contact:vix}{indice de volatilité} à son échéance. Comme l'indice ne peut être ni acheté ni détenu, le
contrat n'est pas relié à l'indice du jour par un \omterm{def:m1:basis-roll-and-delivery:carry}{coût de portage}~: il est
le prix auquel le marché situe l'indice à terme.
\end{definition}

Deux faits structurent ce marché. La volatilité revient vers sa moyenne, de sorte que la
courbe des contrats à terme monte en général lorsque l'indice est bas (\omterm{def:m1:basis-roll-and-delivery:contango}{contango}) et baisse lorsqu'il
a flambé (\omterm{def:m1:basis-roll-and-delivery:contango}{backwardation}). Les vendeurs de volatilité sont en moyenne rémunérés par une
prime~: en période calme, une position vendeuse sur le premier contrat à terme rapporte le
roll le long de la courbe, mois après mois, comme le \omterm{def:m1:basis-roll-and-delivery:contango}{contango} de la
\cref{fig:m1:basis-roll-and-delivery:rolled} vu de l'autre côté.

\begin{omfigure}
\begin{tikzpicture}
\begin{axis}[width=10.5cm,height=4.8cm,
  xlabel={mois avant l'échéance (0 = l'indice)},ylabel={points de volatilité},
  tick label style={font=\small},label style={font=\small},
  xmin=0,xmax=7,ymin=10,ymax=40,grid=major,grid style={gray!25},
  legend columns=2,legend style={font=\footnotesize,at={(0.5,-0.32)},anchor=north,draw=none,fill=none,/tikz/every even column/.append style={column sep=8pt}},
  table/col sep=comma]
\addplot[omDef,very thick,mark=*] table[x=month,y=calm]{figdata/markets-1/25-volatility-first-contact/term.csv};
\addplot[omThm,very thick,dashed,mark=square*] table[x=month,y=stress]{figdata/markets-1/25-volatility-first-contact/term.csv};
\legend{un marché calme~: contango,après un pic~: backwardation}
\end{axis}
\end{tikzpicture}

\omcaption{Deux courbes stylisées de contrats à terme sur un \omterm{def:m1:volatility-first-contact:vix}{indice de volatilité}. Dans la courbe calme, une
position vendeuse sur le premier contrat rapporte 1.6 point par mois si rien ne se passe~; dans l'autre, le
marché anticipe une baisse de l'indice d'un tiers en trois mois.
Chiffres illustratifs.
}\label{fig:m1:volatility-first-contact:term}
\end{omfigure}

\begin{example}[5 février 2018]\label{ex:m1:volatility-first-contact:feb2018}
Le lundi 5 février 2018, le S\&P 500 a chuté de 4~\% et l'\omterm{def:m1:volatility-first-contact:vix}{indice de volatilité} a progressé
de 20 points dans la journée. Les produits cotés en bourse offrant une exposition
à effet de levier ou inverse sur les premiers contrats à terme, avec environ 4~milliards de dollars d'encours à la
fin de 2017, ont dû se rééquilibrer en fin de journée et, selon l'arithmétique de la
\cref{prop:m1:exchange-traded-funds:rebalance}, les deux catégories ont dû \emph{acheter} des contrats à terme après une hausse.
Une analyse d'une banque centrale décrit la
boucle qui en résulte~: leurs achats ont fait monter les contrats à terme, ce qui a accru
le montant qu'ils devaient acheter. Le plus grand produit inverse a perdu presque toute sa valeur
au cours de cette séance, puis son émetteur l'a liquidé.
\end{example}

\section{Tutoriel~: du tableau de cours à un terme à terme, un dividende et un smile}\label{tut:m1:volatility-first-contact:read}

\begin{tutorial}
\textbf{Objectif.} Calculer le prix avec la formule-boîte noire, l'inverser de façon robuste, lire le
terme à terme et le dividende implicite dans la parité, retrouver le smile, et calculer une
volatilité de type indice. \textbf{État final~:} les quatre figures de données de ce
chapitre.
\begin{tutsteps}
\item \textbf{Le solveur.} D'abord les bornes, puis la dichotomie.
\omcode{code/firm/parity/firm_parity.py}{24}{37}{Volatilité implicite, ou une erreur lorsqu'aucune volatilité ne peut expliquer le prix.}\label{lst:m1:volatility-first-contact:solver}
\item \textbf{La parité comme instrument.}
\omcode{code/firm/parity/firm_parity.py}{40}{48}{Le terme à terme à partir d'un seul strike, et le portage qu'il implique.}\label{lst:m1:volatility-first-contact:parity}
\item \textbf{Lecture d'une chaîne} arrondie au centime, telle qu'un écran l'affiche.
\omcode{code/markets-1/25-volatility-first-contact/python/vol_first.py}{33}{42}{Terme à terme, dividende et smile à partir des seuls prix.}\label{lst:m1:volatility-first-contact:read}
\item \textbf{La bande.}
\omcode{code/firm/parity/firm_parity.py}{65}{80}{Variance implicite sans modèle à partir des cotations hors de la monnaie.}\label{lst:m1:volatility-first-contact:strip}
\end{tutsteps}
\textbf{Pour aller plus loin.} Tronquer la bande à 80 et à 120 et mesurer la part de la volatilité
de type indice ainsi perdue~; élargir ensuite le pas des strikes de 1 à 5 et mesurer l'erreur de
discrétisation. Ces deux effets sont de vraies caractéristiques
des indices publiés.
\end{tutorial}

\section{Construction~: contrôleur de parité et solveur de volatilité implicite}\label{bld:m1:volatility-first-contact:parity}

\begin{build}
\bfield{Rôle} Les deux premières fonctions que tous les systèmes d'options de la mini-entreprise
appellent~: quelle volatilité correspond à ce prix, et cette chaîne est-elle cohérente avec un seul
terme à terme~?
\bfield{Interface} \texttt{price(forward, strike, years, rate, vol, right)}~;
\texttt{implied\_vol(premium, forward, strike, years, rate, right)}~;
\texttt{implied\_forward(call, put, strike, years, rate)}~;
\texttt{implied\_dividends\_pv(spot, forward, years, rate)}~;
\texttt{conversion\_edge(\ldots)} et \texttt{reversal\_edge(\ldots)} à partir des cours
acheteurs et vendeurs~; \texttt{variance\_strip(forward, years, rate, quotes)}.
\bfield{Règles} Le solveur lève une erreur pour une prime située hors des bornes
d'absence d'arbitrage et ne renvoie jamais une volatilité négative ou tronquée sans le signaler.
Les conversions et les reversals sont évalués aux prix auxquels on pourrait réellement négocier~:
cours acheteur du call et cours vendeur du put pour l'un, l'inverse pour l'autre, avec le coût
d'emprunt de la \cref{bld:m1:stock-loan-and-short-selling:borrow} côté vente à découvert.
Les séries et les échéances proviennent du conteneur de chaîne
(\cref{bld:m1:option-contracts-and-exchanges:chain}).
\bfield{Tests d'acceptation} \texttt{code/firm/parity/tests/}~: parité du pricer~; allers-retours du solveur
sur les strikes et les volatilités, ainsi que ses refus~; un dividende retrouvé à partir d'un terme
à terme~; conversions et reversals non rentables à l'intérieur des cotations et rentables
en présence d'une erreur introduite volontairement~; la bande retrouve une volatilité plate
à deux centièmes de point près.
\bfield{Pour aller plus loin} Options américaines~: une prime d'exercice anticipé par un arbre
binomial, et la forme en inégalité de la parité.
\end{build}

\begin{omsources}
\item Cboe Global Indices, \emph{Volatility Index Methodology: Cboe
Volatility Index} (formule, sélection des options, interpolation).
\item Cboe Futures Exchange, \emph{VIX Futures: contract specifications}.
\item Bank for International Settlements, «~The equity market turbulence of 5
February: the role of exchange-traded volatility products~», \emph{BIS
Quarterly Review}, mars 2018.
\item F.~Black et M.~Scholes, «~The pricing of options and corporate
liabilities~», \emph{Journal of Political Economy} 81 (1973)~; F.~Black,
«~The pricing of commodity contracts~», \emph{Journal of Financial Economics}
3 (1976).
\end{omsources}

\section{Exercices}

\begin{exercise}[$\star$]\label{exo:m1:volatility-first-contact:1}
Avec $r = 4\%$ et $T = 0.25$, le call 100 vaut 4.06 et le put 100 vaut 3.87.
Calculer le terme à terme implicite. Avec l'action à 100, calculer le portage implicite.
\end{exercise}

\begin{exercise}[$\star$]\label{exo:m1:volatility-first-contact:2}
Utiliser l'approximation \omterm{def:m1:volatility-first-contact:atm}{à la monnaie} pour estimer la \omterm{def:m1:volatility-first-contact:iv}{volatilité implicite} du call 100 de
l'exercice précédent, puis comparer avec les 20.1 donnés par le solveur.
\end{exercise}

\begin{exercise}[$\star$]\label{exo:m1:volatility-first-contact:3}
Un contrat à terme sur l'\omterm{def:m1:volatility-first-contact:vix}{indice de volatilité} est acheté à 16.70 et vendu à 19.25.
Calculer le P\&L de 40 contrats. Combien de ticks cela représente-t-il~?
\end{exercise}

\begin{exercise}[$\star\star$]\label{exo:m1:volatility-first-contact:4}
Un call 90 à trois mois sur une action à 100 (terme à terme 100.20, $r = 4\%$) est coté
10.05. Montrer qu'il n'admet pas de \omterm{def:m1:volatility-first-contact:iv}{volatilité implicite}. Que doit faire un système avec
cette cotation~?
\end{exercise}

\begin{exercise}[$\star\star$]\label{exo:m1:volatility-first-contact:5}
À partir d'une chaîne, on lit un terme à terme de 98.10 à trois mois, avec l'action à 100
et $r = 4\%$. La société ne verse aucun dividende. Calculer le \omterm{def:m1:financing:seclending}{frais d'emprunt} implicite
en taux annuel, puis l'interpréter.
\end{exercise}

\begin{exercise}[$\star\star$]\label{exo:m1:volatility-first-contact:6}
Les rendements logarithmiques quotidiens d'un indice sur cinq jours sont $+0.4\%$, $-1.1\%$, $+0.7\%$,
$-2.0\%$, $+1.5\%$. Donner la \omterm{def:m1:volatility-first-contact:realised}{volatilité réalisée}, annualisée sur 252 jours et
sans retrancher la moyenne. Si la \omterm{def:m1:volatility-first-contact:iv}{volatilité implicite} valait 16, qui avait raison~?
\end{exercise}

\begin{exercise}[$\star\star\star$]\label{exo:m1:volatility-first-contact:7}
\emph{Programmation.} Avec \texttt{read\_chain} sur la chaîne à neuf strikes du chapitre,
calculer le terme à terme implicite moyen, l'étendue des neuf termes à terme strike par strike,
le dividende implicite et le skew 90--110. Arrondir ensuite les prix à cinq centimes au lieu d'un centime
et donner les mêmes résultats~: quels chiffres résistent~?
\end{exercise}

\begin{exercise}[$\star\star\star$]\label{exo:m1:volatility-first-contact:8}
\emph{Trouver la faille.} «~L'\omterm{def:m1:volatility-first-contact:vix}{indice de volatilité} est à 14 et sa moyenne de long terme
est de 19. Il est bon marché~: acheter le contrat à terme de premier rang à 15.6 et attendre le retour
à la moyenne~». Utiliser la \cref{fig:m1:volatility-first-contact:term}.
\end{exercise}

\section{Problème~: lecture d'une chaîne}

\begin{problem}[{Problème du week-end --- tout ce que l'écran sait}]\label{pb:m1:volatility-first-contact:1}
Une action se négocie à 100.00. Des options européennes à trois mois ($T = 0.25$, $r = 4\%$)
affichent les prix milieu suivants~:

\begin{center}\small
\begin{tabular}{lrrrrr}
\toprule
Strike & 90 & 95 & 100 & 105 & 110 \\
\midrule
Call & 11.22 & 7.26 & 4.06 & 1.88 & 0.68 \\
Put & 1.12 & 2.11 & 3.87 & 6.63 & 10.38 \\
\bottomrule
\end{tabular}
\end{center}

\medskip\noindent\textbf{Partie I --- Le terme à terme.}
\begin{enumerate}
\item Calculer le terme à terme implicite à chacun des cinq strikes.
\item Pourquoi diffèrent-ils légèrement~? Lequel est le moins fiable~?
\item Donner leur moyenne.
\item Donner le portage implicite, en dollars.
\item La société doit verser 0.80 dans cinq semaines. L'action est-elle
difficile à emprunter~?
\end{enumerate}

\medskip\noindent\textbf{Partie II --- Le sourire.}
\begin{enumerate}[resume]
\item Quelle option utilise-t-on à chaque strike pour calculer la \omterm{def:m1:volatility-first-contact:iv}{volatilité implicite},
et pourquoi~?
\item Le solveur donne 23.8, 21.8, 20.1, 18.7 et 17.6. Donner le skew 90--110
et la pente en points de volatilité par 10~\% de strike.
\item Vérifier la valeur \omterm{def:m1:volatility-first-contact:atm}{à la monnaie} avec l'approximation de
la \cref{def:m1:volatility-first-contact:atm}.
\item Un collègue calcule la \omterm{def:m1:volatility-first-contact:iv}{volatilité implicite} du \emph{call} 90 avec le spot
à la place du terme à terme et trouve 25.1 au lieu de 23.8. Expliquer.
\end{enumerate}

\medskip\noindent\textbf{Partie III --- Passer les ordres.}
\begin{enumerate}[resume]
\item Le call 100~: cours acheteur 4.03, cours vendeur 4.09~; le put 100~: 3.84 à 3.90~;
l'action 99.99 à 100.01. Avec le dividende de 0.80, valant 0.797 aujourd'hui, évaluer la
conversion.
\item Évaluer le reversal, avec un coût d'emprunt de 0.05.
\item Que disent les deux résultats sur le marché~?
\item Une cotation apparaît~: le call 100 a un cours acheteur de 4.60. Que faire, et que vérifier
d'abord~?
\item Le lendemain, l'action vaut 97 et le call 100 vaut 2.74. Le call est-il
devenu «~moins cher~»~? Que calculer pour répondre~?
\end{enumerate}

\medskip\noindent\textbf{Partie IV --- Jugement.}
\begin{enumerate}[resume]
\item Pourquoi les desks cotent-ils la volatilité et non le prix~?
\item Que dit un skew de six points sur la vision du marché concernant
l'action~?
\item Sur cette chaîne, une volatilité de type indiciel serait de 23.0 contre 20.1 \omterm{def:m1:volatility-first-contact:atm}{à
la monnaie}. Pourquoi est-elle plus élevée~?
\item Pourquoi la \omterm{def:m1:volatility-first-contact:iv}{volatilité implicite} peut-elle être en moyenne supérieure à la \omterm{def:m1:volatility-first-contact:realised}{volatilité réalisée}
sans que personne ne soit irrationnel~?
\item Énoncer le \emph{résultat nommé}~: le dividende implicite et le
skew 90--110.
\item En une phrase~: qu'est-ce que la \omterm{def:m1:volatility-first-contact:iv}{volatilité implicite}~?
\end{enumerate}
\end{problem}

\section{Questions d'entretien}

\begin{interviewq}[$\star$ \iqroles{trader, chercheur}]\label{iq:m1:volatility-first-contact:1}
Énoncez la \omterm{prop:m1:volatility-first-contact:parity}{parité call-put} et dites ce qu'elle suppose.
\end{interviewq}

\begin{interviewq}[$\star$ \iqroles{trader, chercheur}]\label{iq:m1:volatility-first-contact:2}
Qu'est-ce que la \omterm{def:m1:volatility-first-contact:iv}{volatilité implicite}, et pourquoi n'est-elle pas une prévision~?
\end{interviewq}

\begin{interviewq}[$\star\star$ \iqroles{développeur, chercheur}]\label{iq:m1:volatility-first-contact:3}
Écrivez un solveur de \omterm{def:m1:volatility-first-contact:iv}{volatilité implicite}. Qu'est-ce qui peut mal se passer~?
\end{interviewq}

\begin{interviewq}[$\star\star$ \iqroles{trader, chercheur}]\label{iq:m1:volatility-first-contact:4}
Les calls paraissent bon marché par rapport aux puts au même strike. Quelles sont les explications
possibles avant de vous prononcer sur un arbitrage~?
\end{interviewq}

\begin{interviewq}[$\star\star$ \iqroles{chercheur, trader}]\label{iq:m1:volatility-first-contact:5}
Comment un \omterm{def:m1:volatility-first-contact:vix}{indice de volatilité} est-il calculé, et pourquoi ne pouvez-vous pas l'acheter~?
\end{interviewq}

\begin{interviewq}[$\star\star\star$ \iqroles{trader, chercheur}]\label{iq:m1:volatility-first-contact:6}
Expliquez ce qui s'est passé avec les produits de vente de volatilité en février 2018, et
ce qu'un gestionnaire des risques aurait dû voir à l'avance.
\end{interviewq}
